\section{A Router With Nothing to Serve}
\label{sec:ch08-a-router-with-nothing-to-serve}

\index{Cosmo Router!where it comes from|(}%
The router is not a package. Chapter~\ref{ch:composition} added \code{wgc}
through npm and chapter~\ref{ch:first-subgraph} added a NuGet reference, and
this chapter adds neither: the Cosmo Router is a compiled binary published as a
GitHub release asset, and you download it. It is also published as a container
image, and this book does not use that one. Seven platform and architecture
pairs ship as archives with each release, macOS, Linux and Windows on the
architectures you would expect, so nothing in this chapter needs Docker
installed.

\index{Cosmo Router!the archive holds two files}%
The Windows archive holds exactly two files, and one of them is the Apache
license. The other is the router. The release publishes an md5 beside each
asset, so check it before you unzip; the version table in
appendix~\ref{app:toolchain} pins which release this book was written against.

\index{Cosmo Router!what the binary reports about itself}%
Ask the binary what it is:

\begin{minted}{text}
router -version
\end{minted}

\begin{minted}{text}
Router:
  Version: 0.341.0
  Go version: go1.26.6
  OS: windows
  Arch: amd64
  Built: 2026-08-18T12:27:33Z
  VCS Revision: e7b12b1beac4cde7f80cca8b12fdd250ef50cc63
  Dependencies:
    github.com/wundergraph/astjson v1.1.0
    github.com/wundergraph/graphql-go-tools/v2 v2.16.0
\end{minted}

\index{Cosmo Router!names its own query engine}%
Two dependencies, and the second one is the query engine. The planner this
chapter spends its second half reading is in \code{graphql-go-tools}, which is
a separate library on its own version, along with the normalizer and the thing
that answers introspection. Worth noting before you go looking for any of the
three in the router's own source tree and do not find them.

\subsection{It Gets Most of the Way, Then Stops}

\index{Cosmo Router!refuses to start with no config}%
Run it in an empty directory with no arguments at all. It does not fail
immediately, which is the interesting part:

\begin{minted}[fontsize=\footnotesize,breakanywhere]{text}
INFO  Config file watching is disabled, you can still trigger reloads by sending SIGHUP to the router process
WARN  GOMEMLIMIT was not set. Please set it manually to around 90% of the available memory to prevent OOM kills
      error="failed to set GOMEMLIMIT: cgroups is not supported on this system"
WARN  No graph token provided. The following Cosmo Cloud features are disabled. Not recommended for Production.
      features=["Schema Usage Tracking","Persistent operations","Advanced Request Tracing","Cosmo Cloud Tracing","Cosmo Cloud Metrics"]
WARN  Net poller is only available on Linux and MacOS. Falling back to less efficient connection handling method.
      error="epoll/kqueue is not supported on this system"
INFO  Prometheus metrics enabled
      listen_addr="127.0.0.1:8088" endpoint="/metrics"
INFO  Serving GraphQL playground
      url="http://localhost:3002/"
ERROR Could not start router
      error="router start error: failed to start router: failed to bootstrap router: graph token is required to fetch execution config from CDN. Alternatively, configure a custom storage provider or specify a static execution config"
\end{minted}

\index{Cosmo Router!log lines are newline-delimited json}%
Those lines arrive as newline-delimited JSON, one object per line, and every
one of them also carries a hostname, a pid, a service name and a service
version. The block above is those lines with exactly those four fields removed
and everything else indented underneath, which is the only change made to it
and the same change made to every router log this chapter prints.

\index{Cosmo Router!its default source is a CDN}%
Read the last line first. \enquote{Graph token is required to fetch execution
config from CDN.} A router given nothing does not look for a file. It goes to
WunderGraph's CDN, and a graph token is what buys an answer there. The default
configuration of this binary is a client of the vendor's platform, and running
it locally is the deviation rather than the baseline. That is a reasonable
default for a company selling a control plane and it is not the mode this book
is in, so meet the refusal here rather than under some later error that does
not explain itself this well.

\index{Cosmo Router!names the local alternative in the refusal}%
Then read the rest of the same sentence, because it is unusually good error
text. \enquote{Alternatively, configure a custom storage provider or specify a
static execution config.} The local path is named in the refusal. This is not a
tool being worked around, and the next section is one file long.

\subsection{Three Things It Decided Before It Failed}

\index{Cosmo Router!announces its port before failing}%
\index{ports!the router takes its own default}%
It announced \code{http://localhost:3002/} on the way down. That is the
router's own default address, chosen with no configuration present at all, and
it is the address this book uses. Chapter~\ref{ch:hot-chocolate-zero} had to pin
5001 in a launch profile because \code{dotnet new web} picks a port at random
and a book printing raw HTTP would otherwise print fiction. The router has no
such problem: it prints the same address on every machine, so pinning one in
\code{config.yaml} would add a line that can go stale against the log beside it
and buy nothing.

\index{Cosmo Router!serves a playground at the root}%
That address is the playground, at \code{/}. GraphQL is at \code{/graphql},
which is the path chapter~\ref{ch:hot-chocolate-zero} put the service on, so
the two differ by port and by nothing else.

\index{Cosmo Router!Prometheus metrics are on by default}%
And it opened a second listener, for Prometheus metrics, on loopback port 8088
at \code{/metrics}. It did that before it had a graph to serve and without
being asked. That is the third time in three chapters a library has written
something into a public surface on its own initiative:
chapter~\ref{ch:first-subgraph} found \code{@shareable} appearing on a generated
type, chapter~\ref{ch:composition} found \code{@cost} on every field with a
resolver behind it, and here is a metrics endpoint on a port nobody chose.
Name it now, because chapter~\ref{ch:operations} has to decide what a subgraph
behind a router should expose, and the honest starting position is that
something already is.

\index{Advanced Request Tracing!is in the disabled list}%
One more line matters later. Five Cosmo Cloud features are named as disabled
without a token, and the third of them is Advanced Request Tracing. Nothing in
this chapter needs the other four. That one costs us the query plan, and
section~\ref{sec:ch08-the-plan-with-one-fetch-in-it} is where the bill arrives.
\index{Cosmo Router!where it comes from|)}%
